% ============================================================
%  CALCULUS – Typographic Design
%  Compile with: pdflatex -shell-escape calculus.tex (2x)
% ============================================================
\documentclass[a4paper,10pt,twoside]{book}

%% ─── Packages ───────────────────────────────────────────────
\usepackage[T1]{fontenc}
\usepackage[utf8]{inputenc}

\usepackage{fix-cm}
\usepackage{helvet}
\renewcommand{\familydefault}{\sfdefault}
\usepackage[top=2.2cm,bottom=2.5cm,inner=2.4cm,outer=1.8cm,
            marginparwidth=0cm,includeheadfoot,headheight=21pt]{geometry}
\usepackage{xcolor}
\usepackage{tikz}
\usetikzlibrary{calc,positioning,shapes.geometric,decorations.pathmorphing}
\usepackage{pgfplots}
\pgfplotsset{compat=1.18}
\usepackage{tcolorbox}
\tcbuselibrary{skins,breakable,raster,hooks}
\usepackage{fancyhdr}
\usepackage{titlesec}
\usepackage{titletoc}
\usepackage{amsmath,amssymb,amsthm}
\usepackage{mathtools}
\usepackage{enumitem}
\usepackage{multicol}
\usepackage{wrapfig}
\usepackage{adjustbox}
\usepackage{pifont}
\usepackage{marvosym}
\usepackage{tasks}
\usepackage{qrcode}
\usepackage{microtype}
\usepackage{graphicx}
\usepackage{booktabs}
\usepackage{colortbl}
\usepackage{array}
\usepackage{tabularx}
\usepackage{varwidth}
\setlength{\parindent}{0pt}
\setlength{\parskip}{4pt}

%% ─── Color Palette ──────────────────────────────────────────
\definecolor{Cap1Blue}{HTML}{1E9ED7}   % turquoise blue
\definecolor{Cap2Orange}{HTML}{F07E26} % orange
\definecolor{Cap3Green}{HTML}{5BB85D}  % green
\definecolor{Cap4Purple}{HTML}{7C5CBF} % purple
\definecolor{AccentYellow}{HTML}{FFD93D}
\definecolor{AccentRed}{HTML}{E84855}
\definecolor{LightBlue}{HTML}{DAEEF8}
\definecolor{LightOrange}{HTML}{FDEBD0}
\definecolor{LightGreen}{HTML}{D5F5D5}
\definecolor{LightPurple}{HTML}{EAE4F7}
\definecolor{LightYellow}{HTML}{FFFBCC}
\definecolor{DarkText}{HTML}{1A1A2E}
\definecolor{MidGray}{HTML}{6C757D}
\definecolor{BoxGray}{HTML}{F4F6F8}
\definecolor{RuleGray}{HTML}{DEE2E6}

%% ─── Chapter Color (redefined per chapter) ──────────────────
\newcommand{\ChapColor}{Cap1Blue}
\newcommand{\ChapColorLight}{LightBlue}
\newcommand{\setChapColor}[2]{\renewcommand{\ChapColor}{#1}%
                               \renewcommand{\ChapColorLight}{#2}}

%% ─── Header / Footer ────────────────────────────────────────
\fancypagestyle{main}{%
  \fancyhf{}
  \fancyhead[LE]{%
    \begin{tikzpicture}[baseline=-0.5ex]
      \fill[\ChapColor] (0,0) rectangle (0.5cm,0.45cm);
      \node[anchor=west,text=white,font=\bfseries\small] at (0.55cm,0.22cm)
        {\nouppercase{\leftmark}};
    \end{tikzpicture}}
  \fancyhead[RO]{%
    \begin{tikzpicture}[baseline=-0.5ex]
      \node[anchor=east,text=MidGray,font=\small] at (-0.55cm,0.22cm)
        {\nouppercase{\rightmark}};
      \fill[\ChapColor] (0,0) rectangle (-0.5cm,0.45cm);
    \end{tikzpicture}}
  \fancyfoot[LE]{%
    \begin{tikzpicture}[baseline=-0.5ex]
      \fill[\ChapColor,rounded corners=10pt] (0,0) rectangle (0.65cm,0.5cm);
      \node[text=white,font=\bfseries\small] at (0.325cm,0.25cm) {\thepage};
    \end{tikzpicture}}
  \fancyfoot[RO]{%
    \begin{tikzpicture}[baseline=-0.5ex]
      \fill[\ChapColor,rounded corners=10pt] (0,0) rectangle (0.65cm,0.5cm);
      \node[text=white,font=\bfseries\small] at (0.325cm,0.25cm) {\thepage};
    \end{tikzpicture}}
  \fancyfoot[C]{\color{MidGray}\tiny Calculus}
  \renewcommand{\headrulewidth}{0pt}
  \renewcommand{\footrulewidth}{0pt}}
\pagestyle{main}

%% ─── Chapter Style ──────────────────────────────────────────
\titleformat{\chapter}[block]{}{}
  {0pt}{\chaptertitleformat}
\newcommand{\chaptertitleformat}[1]{%
  \begin{tikzpicture}[remember picture,overlay]
    %% lateral stripe
    \fill[\ChapColor] (-1cm,0.2cm) rectangle (0.1cm,-5.5cm);
    %% large circle with number
    \fill[\ChapColor] (1.2cm,-1.8cm) circle (1.5cm);
    \fill[white,opacity=0.15] (1.2cm,-1.8cm) circle (1.2cm);
    \node[text=white,font=\fontsize{52}{52}\selectfont\bfseries] at (1.2cm,-1.8cm)
      {\thechapter};
    %% title in uppercase
    \node[anchor=north west,text=\ChapColor,font=\fontsize{28}{32}\selectfont\bfseries,
          align=left,text width=10cm] at (3.2cm,-0.3cm) {#1};
    %% decorative wide stripe
    \fill[\ChapColor,opacity=0.12] (-1cm,-4.5cm) rectangle (\textwidth,-5.5cm);
  \end{tikzpicture}%
  \vspace{5.8cm}}

\titlespacing{\chapter}{0pt}{0pt}{0pt}

\titleformat{\section}[block]{\large\bfseries\color{\ChapColor}}
  {\makebox[1.2cm][l]{\thesection}}{0pt}{}[{%
    \color{\ChapColor}\rule[0.4ex]{\linewidth}{1.5pt}}]

\titleformat{\subsection}[runin]{\bfseries\color{DarkText}}
  {\thesubsection\;}{0pt}{}[~---]

%% ─── tcolorbox Environments ─────────────────────────────────

%% Definition
\tcbset{definicaobase/.style={
    enhanced,breakable,
    colback=LightBlue,colframe=Cap1Blue,
    colbacktitle=Cap1Blue,coltitle=white,
    fonttitle=\bfseries\small,
    arc=6pt,boxrule=1.2pt,
    attach boxed title to top left={xshift=8pt,yshift=-3pt},
    boxed title style={arc=4pt,boxrule=0pt},
    left=6pt,right=6pt,top=10pt,bottom=6pt}}

\newenvironment{definition}[1]{%
  \begin{tcolorbox}[definicaobase,
    colback=\ChapColorLight,colframe=\ChapColor,
    colbacktitle=\ChapColor,
    title={\ding{100}\ Definition\ifx&#1&\else\ ---\ #1\fi}]}
  {\end{tcolorbox}}

%% Theorem
\tcbset{teoremabase/.style={
    enhanced,breakable,
    colback=LightYellow,colframe=AccentYellow,
    colbacktitle=AccentYellow,coltitle=DarkText,
    fonttitle=\bfseries\small,
    arc=6pt,boxrule=1.2pt,
    attach boxed title to top left={xshift=8pt,yshift=-3pt},
    boxed title style={arc=4pt,boxrule=0pt},
    left=6pt,right=6pt,top=10pt,bottom=6pt}}

\newenvironment{theorem}[1]{%
  \begin{tcolorbox}[teoremabase,title={\ding{72}\ Theorem\ifx&#1&\else\ ---\ #1\fi}]}
  {\end{tcolorbox}}

%% Example
\newenvironment{example}[1]{%
  \begin{tcolorbox}[enhanced,breakable,
    colback=BoxGray,colframe=RuleGray,
    colbacktitle=MidGray,coltitle=white,
    fonttitle=\bfseries\small,
    arc=6pt,boxrule=0.6pt,
    attach boxed title to top left={xshift=8pt,yshift=-3pt},
    boxed title style={arc=4pt,boxrule=0pt},
    left=6pt,right=6pt,top=10pt,bottom=6pt,
    title={\ding{228}\ Example\ifx&#1&\else\ #1\fi}]}
  {\end{tcolorbox}}

%% Caution / Warning
\newenvironment{caution}[0]{%
  \begin{tcolorbox}[enhanced,
    colback=LightOrange,colframe=Cap2Orange,
    colbacktitle=Cap2Orange,coltitle=white,
    fonttitle=\bfseries\small,
    arc=6pt,boxrule=1.2pt,
    attach boxed title to top left={xshift=8pt,yshift=-3pt},
    boxed title style={arc=4pt,boxrule=0pt},
    left=6pt,right=6pt,top=10pt,bottom=6pt,
    title={\ding{45}\ Caution!}]}
  {\end{tcolorbox}}

%% Numbered Exercise
\newcounter{exerc}[chapter]
\renewcommand{\theexerc}{\thechapter.\arabic{exerc}}

\newenvironment{exercise}[1][]{%
  \refstepcounter{exerc}%
  \begin{tcolorbox}[enhanced,breakable,
    colback=white,colframe=\ChapColor,
    arc=6pt,boxrule=1pt,
    left=6pt,right=6pt,top=6pt,bottom=6pt,
    before upper={\textbf{\color{\ChapColor}\theexerc}\ifx&#1&\else\ ---\ \textit{#1}\fi\par\smallskip}]}
  {\end{tcolorbox}}

%% Did You Know / Further Reading
\newenvironment{didyouknow}[0]{%
  \begin{tcolorbox}[enhanced,
    colback=LightGreen,colframe=Cap3Green,
    colbacktitle=Cap3Green,coltitle=white,
    fonttitle=\bfseries\small,
    arc=10pt,boxrule=1.2pt,
    attach boxed title to top left={xshift=8pt,yshift=-3pt},
    boxed title style={arc=6pt,boxrule=0pt},
    left=6pt,right=6pt,top=10pt,bottom=6pt,
    title={\ding{108}\ Did You Know?}]}
  {\end{tcolorbox}}

%% Pro / Con Comparison Box
\newcommand{\procon}[4]{%
  \begin{tcolorbox}[enhanced,sidebyside,sidebyside align=top,
    colback=white,colframe=RuleGray,arc=6pt,boxrule=0.6pt,
    lefthand width=0.48\linewidth]
    \textbf{\color{Cap3Green}\ding{52}\ #1}\\[4pt]#2
    \tcblower
    \textbf{\color{AccentRed}\ding{55}\ #3}\\[4pt]#4
  \end{tcolorbox}}

%% Highlighted Numerical Result
\newcommand{\result}[1]{%
  \begin{center}
  \begin{tcolorbox}[enhanced,width=0.55\linewidth,
    colback=\ChapColorLight,colframe=\ChapColor,
    arc=8pt,boxrule=1.5pt,halign=center,
    left=8pt,right=8pt,top=6pt,bottom=6pt]
    \large\bfseries\color{\ChapColor}#1
  \end{tcolorbox}
  \end{center}}

%% Numbered Step Block
\newcommand{\step}[2]{%
  \begin{tikzpicture}[baseline=(n.base)]
    \fill[\ChapColor] (0,0) circle (0.25cm);
    \node[text=white,font=\bfseries\small] (n) at (0,0) {#1};
  \end{tikzpicture}\hspace{6pt}#2\par\smallskip}

%% Highlighted Formula Box (orange)
\newcommand{\formula}[2][]{%
  \begin{tcolorbox}[enhanced,
    colback=LightOrange,colframe=Cap2Orange,
    arc=6pt,boxrule=1.5pt,halign=center,
    left=8pt,right=8pt,top=4pt,bottom=4pt]
    \ifx&#1&\else\textbf{\small\color{Cap2Orange}#1}\\[4pt]\fi
    \( \displaystyle #2 \)
  \end{tcolorbox}}

%% QR Code with Caption
\newcommand{\qrbox}[2]{%
  \begin{tcolorbox}[enhanced,width=3.8cm,
    colback=white,colframe=\ChapColor,
    arc=6pt,boxrule=1pt,halign=center,
    left=4pt,right=4pt,top=4pt,bottom=4pt]
    \qrcode[height=2.5cm]{#1}\\[4pt]
    {\tiny\color{MidGray}#2}
  \end{tcolorbox}}

%% Lesson Box
\newcounter{lesson}[chapter]
\renewcommand{\thelesson}{Lesson\ \thechapter.\arabic{lesson}}
\newenvironment{lesson}[1]{%
  \refstepcounter{lesson}%
  \begin{tcolorbox}[enhanced,breakable,
    colback=white,colframe=\ChapColor,
    colbacktitle=\ChapColor,coltitle=white,
    fonttitle=\bfseries\small,
    arc=8pt,boxrule=1.5pt,
    attach boxed title to top left={xshift=10pt,yshift=-4pt},
    boxed title style={arc=5pt,boxrule=0pt},
    left=6pt,right=6pt,top=12pt,bottom=6pt,
    title={\thelesson\ ---\ #1}]}
  {\end{tcolorbox}}

%% Quick Answer Key
\newenvironment{answerkey}{%
  \smallskip
  \begin{tcolorbox}[enhanced,
    colback=\ChapColorLight,colframe=\ChapColor,
    arc=6pt,boxrule=0.8pt,
    title={\small\bfseries\color{white} Quick Answer Key},
    colbacktitle=\ChapColor,fonttitle=\small,
    left=4pt,right=4pt,top=6pt,bottom=4pt,
    attach boxed title to top left={xshift=8pt,yshift=-3pt},
    boxed title style={arc=4pt}]
  \footnotesize\begin{multicols}{3}}
  {\end{multicols}\end{tcolorbox}}

%% ─── Styled TOC ─────────────────────────────────────────────
\titlecontents{chapter}[0pt]
  {\vspace{6pt}\bfseries}
  {\begin{tikzpicture}[baseline=-0.5ex,inner sep=0pt]
     \fill[Cap1Blue] (0,0) circle (0.22cm);
     \node[text=white,font=\tiny\bfseries] at (0,0) {\thecontentslabel};
   \end{tikzpicture}\hspace{6pt}}
  {}{\titlerule*[6pt]{.}\contentspage}

\titlecontents{section}[1.4cm]
  {}
  {\contentslabel{1.2cm}}
  {}{\titlerule*[6pt]{.}\contentspage}

%% ─── Decorative Section Divider ─────────────────────────────
\newcommand{\secbar}{%
  \begin{tikzpicture}
    \fill[\ChapColor] (0,0) rectangle (\linewidth,0.08cm);
    \fill[AccentYellow] (0,0) rectangle (2cm,0.08cm);
  \end{tikzpicture}\par\smallskip}

%% ─── Topic Icon ─────────────────────────────────────────────
\newcommand{\topic}[2]{%
  \begin{tikzpicture}[baseline=-0.5ex]
    \fill[\ChapColor,rounded corners=3pt] (0,0) rectangle (0.35cm,0.35cm);
    \node[text=white,font=\bfseries\tiny] at (0.175cm,0.175cm) {#1};
  \end{tikzpicture}\hspace{4pt}\textbf{#2}}

%% ─── Styled Lists ───────────────────────────────────────────
\setlist[itemize,1]{label=\textcolor{\ChapColor}{\ding{228}},
                    leftmargin=1.2em,itemsep=2pt}
\setlist[enumerate,1]{label=\textcolor{\ChapColor}{\bfseries\arabic*.},
                      leftmargin=1.5em,itemsep=2pt}

%% ─── End of Preamble ────────────────────────────────────────
\begin{document}

%%=============================================================
%%  COVER PAGE
%%=============================================================
\begin{titlepage}
\begin{tikzpicture}[remember picture,overlay]
  %% Background
  \fill[Cap1Blue] (current page.north west) rectangle (current page.south east);
  %% Decorative shapes
  \fill[Cap1Blue!70!black] (current page.north east) +(-3cm,-6cm)
    [rounded corners=2cm] rectangle +(0,0);
  \fill[white,opacity=0.08] (current page.south west) +(4cm,5cm) circle (4.5cm);
  \fill[white,opacity=0.06] (current page.north east) +(-5cm,-10cm) circle (6cm);
  \fill[Cap2Orange] (current page.west) +(0,-2cm)
    rectangle +(0.7cm,\paperheight);
  \fill[AccentYellow] (current page.west) +(0.7cm,-2cm)
    rectangle +(1.1cm,\paperheight);
  %% Large decorative circle
  \draw[white,opacity=0.15,line width=12pt]
    (current page.center) +(3cm,2cm) circle (6cm);
  \draw[white,opacity=0.08,line width=20pt]
    (current page.center) +(3cm,2cm) circle (8cm);
  %% Title
  \node[anchor=north west,text=white,font=\fontsize{60}{64}\selectfont\bfseries,
        align=left] at ($(current page.north west)+(2.0cm,-3.5cm)$)
    {CALCULUS};
  \node[anchor=north west,text=AccentYellow,
        font=\fontsize{22}{26}\selectfont,align=left]
    at ($(current page.north west)+(2.1cm,-6.4cm)$)
    {A modern course with\\applications and visualizations};
  %% Decorative line
  \draw[white,line width=1.5pt]
    ($(current page.north west)+(2.0cm,-5.8cm)$)
    -- ($(current page.north west)+(14cm,-5.8cm)$);
  %% Chapter summary icons
  \foreach \n/\cor/\nome/\y in {
    1/Cap1Blue/Limits and Continuity/-9.2cm,
    2/Cap2Orange/Derivatives/-10.6cm,
    3/Cap3Green/Applications of Derivatives/-12.0cm,
    4/Cap4Purple/Indefinite Integrals/-13.4cm}{
    \fill[\cor] ($(current page.north west)+(2.1cm,\y)$) circle (0.4cm);
    \node[text=white,font=\bfseries] at ($(current page.north west)+(2.1cm,\y)$) {\n};
    \node[anchor=west,text=white,font=\large]
      at ($(current page.north west)+(2.75cm,\y)$) {\nome};
  }
  %% Decorative graph
  \begin{scope}[shift={($(current page.south east)+(-5.5cm,4cm)$)}]
    \begin{axis}[
      width=7cm,height=6cm,
      axis line style={white,opacity=0.4},
      tick style={white,opacity=0.3},
      ticklabel style={white,opacity=0.4,font=\tiny},
      domain=-3.5:3.5,samples=80,
      xmin=-3.8,xmax=3.8,ymin=-1.5,ymax=1.5,
      grid=none]
      \addplot[white,line width=1.5pt,opacity=0.7] {sin(deg(x))};
      \addplot[AccentYellow,line width=1.5pt,opacity=0.7] {cos(deg(x))};
      \addplot[Cap2Orange,line width=1.5pt,opacity=0.7,domain=-2:2] {x^2/4-1};
    \end{axis}
  \end{scope}
  %% Footer
  \node[anchor=south west,text=white,font=\small,opacity=0.7]
    at ($(current page.south west)+(2.0cm,0.8cm)$) {Calculus for Science and Engineering\ \textbullet\ 1st edition};
\end{tikzpicture}
\end{titlepage}

%%=============================================================
%%  TABLE OF CONTENTS
%%=============================================================
\frontmatter
\tableofcontents

%%=============================================================
%%  CHAPTER 1 — LIMITS AND CONTINUITY
%%=============================================================
\mainmatter
\setChapColor{Cap1Blue}{LightBlue}
\chapter{Limits and Continuity}
\markboth{Ch.~1\ --\ Limits and Continuity}{}

%%─── Chapter Opening ─────────────────────────────────────────
\begin{multicols}{2}
\textbf{In this chapter} you will study the fundamental concept of a \emph{limit},
which underpins all of differential and integral analysis.
The main topics are:

\begin{itemize}
  \item \topic{1}{Intuitive idea of a limit}
  \item \topic{2}{Formal definition ($\varepsilon$-$\delta$)}
  \item \topic{3}{One-sided limits}
  \item \topic{4}{Limits at infinity}
  \item \topic{5}{Continuity}
\end{itemize}

\columnbreak
\begin{didyouknow}
Why are limits important? The derivative, the integral, and the sum of a series
are all defined through a \textbf{limit} process.
Without this concept, Calculus would not exist!
\end{didyouknow}
\end{multicols}

\secbar

%%─── Lesson 1.1 ──────────────────────────────────────────────
\begin{lesson}{Intuitive Idea of a Limit}

The \textbf{limit} of a function $f(x)$ as $x$ approaches $a$
describes the behavior of $f$ \emph{near} $a$, regardless
of what happens \emph{at} $a$.

\medskip
\step{1}{Analyze the table of values below for $f(x)=\frac{x^2-1}{x-1}$:}

\begin{center}
\renewcommand{\arraystretch}{1.3}
\begin{tabular}{|>{\columncolor{LightBlue}}c|c||c|>{\columncolor{LightBlue}}c|}
\hline
\rowcolor{Cap1Blue}
\color{white}\textbf{$x$} & \color{white}\textbf{$f(x)$} &
\color{white}\textbf{$x$} & \color{white}\textbf{$f(x)$} \\
\hline
0.9   & 1.9   & 1.1   & 2.1   \\
0.99  & 1.99  & 1.01  & 2.01  \\
0.999 & 1.999 & 1.001 & 2.001 \\
\hline
\end{tabular}
\end{center}

\step{2}{Notice that the values approach \textbf{2} as $x\to 1$:}

\formula[Limit Notation]{%
  \lim_{x \to 1} \frac{x^2-1}{x-1} = 2}

\end{lesson}

%%─── Formal Section ──────────────────────────────────────────
\section{Formal Definition of a Limit}
\label{sec:lim-formal}

\begin{definition}{Limit ($\varepsilon$-$\delta$ Definition)}
Let $f$ be defined on a neighborhood of $a$ (except possibly at $a$).
We say that
\[
  \lim_{x \to a} f(x) = L
\]
if, for every $\varepsilon > 0$, there exists $\delta > 0$ such that
\[
  0 < |x - a| < \delta \implies |f(x) - L| < \varepsilon.
\]
\end{definition}

\begin{multicols}{2}
\begin{example}{1.1}
Prove that $\lim_{x\to 3}(2x-1)=5$.

\smallskip\textbf{Solution:}\\
Given $\varepsilon>0$, take $\delta = \varepsilon/2$.
If $0<|x-3|<\delta$, then
\[
  |(2x-1)-5| = 2|x-3| < 2\delta = \varepsilon.
\]
\end{example}

\columnbreak
\begin{caution}
The $\varepsilon$-$\delta$ definition requires the limit to exist
on both sides of $a$. If the one-sided limits
$\lim_{x\to a^-}f(x)$ and $\lim_{x\to a^+}f(x)$
differ, the limit \textbf{does not exist}.
\end{caution}
\end{multicols}

%%─── Inline graph with pgfplots ──────────────────────────────
\begin{center}
\begin{tikzpicture}
\begin{axis}[
  width=0.72\linewidth, height=5.5cm,
  xlabel={$x$}, ylabel={$f(x)$},
  xlabel style={font=\small},
  ylabel style={font=\small},
  tick label style={font=\small},
  axis lines=center,
  xmin=-1.5, xmax=3.5,
  ymin=-0.5, ymax=4,
  samples=120,
  grid=both,
  grid style={RuleGray,line width=0.3pt},
  legend style={at={(0.98,0.05)},anchor=south east,font=\small,
                fill=white,draw=RuleGray}]
  \addplot[Cap1Blue,line width=1.8pt,domain=-1.4:2.99] {x+1};
  \addplot[Cap1Blue,line width=1.8pt,domain=3.01:3.4] {x+1};
  \addplot[Cap2Orange,only marks,mark=*,mark size=2.5pt]
    coordinates {(1,2)};
  \addplot[white,only marks,mark=*,mark size=1.8pt,
           mark options={draw=Cap1Blue,fill=white}]
    coordinates {(3,4)};
  \addplot[Cap1Blue,only marks,mark=*,mark size=2.5pt]
    coordinates {(3,2)};
  \legend{$f(x)=x+1$ (with removable discontinuity)}
\end{axis}
\end{tikzpicture}
\par\smallskip\textit{
Example of a function with a removable discontinuity at $x=3$.}
\end{center}

\secbar

%%─── Limit Theorems ──────────────────────────────────────────
\section{Properties of Limits}

\begin{theorem}{Algebra of Limits}
If $\lim_{x\to a}f(x)=L$ and $\lim_{x\to a}g(x)=M$, then:
\begin{multicols}{2}
\begin{enumerate}
  \item $\lim_{x\to a}[f(x)+g(x)] = L+M$
  \item $\lim_{x\to a}[f(x)\cdot g(x)] = L\cdot M$
  \item $\lim_{x\to a}\dfrac{f(x)}{g(x)} = \dfrac{L}{M}$,\ \ $M\neq 0$
  \item $\lim_{x\to a}[c\,f(x)] = cL$
\end{enumerate}
\end{multicols}
\end{theorem}

\begin{theorem}{Squeeze Theorem (Sandwich Theorem)}
If $g(x)\le f(x)\le h(x)$ on a neighborhood of $a$, and
$\lim_{x\to a}g(x)=\lim_{x\to a}h(x)=L$, then
$\lim_{x\to a}f(x)=L$.
\end{theorem}

\begin{example}{1.2 — Fundamental Trigonometric Limit}
Show that $\displaystyle\lim_{x\to 0}\frac{\sin x}{x}=1$.

\smallskip
\textbf{Idea:} Use the Squeeze Theorem with
$\cos x \le \dfrac{\sin x}{x} \le 1$ for $0<|x|<\pi/2$.

\result{$\displaystyle\lim_{x\to 0}\dfrac{\sin x}{x}=1$}
\end{example}

\secbar

%%─── Continuity ──────────────────────────────────────────────
\section{Continuity}

\begin{definition}{Continuity at a Point}
A function $f$ is \textbf{continuous at $a$} if:
\begin{enumerate}
  \item $f(a)$ is defined;
  \item $\lim_{x\to a}f(x)$ exists;
  \item $\lim_{x\to a}f(x) = f(a)$.
\end{enumerate}
\end{definition}

\procon{Continuous Function}{%
  Graph with no ``jumps''\\
  Can compute $f(a)$ directly\\
  Theorems like IVT apply%
}{Discontinuous Function}{%
  Graph with ``holes'', jumps or asymptotes\\
  Limit may not exist\\
  Requires careful analysis%
}

%%─── Exercises ───────────────────────────────────────────────
\section*{Chapter 1 Exercises}
\addcontentsline{toc}{section}{Exercises}

\begin{exercise}[Direct computation]
Evaluate the limits:
\begin{tasks}(2)
  \task $\displaystyle\lim_{x\to 2}(3x^2-x+1)$
  \task $\displaystyle\lim_{x\to 0}\frac{\tan x}{x}$
  \task $\displaystyle\lim_{x\to\infty}\frac{2x^2-1}{x^2+3}$
  \task $\displaystyle\lim_{x\to 4}\frac{\sqrt{x}-2}{x-4}$
\end{tasks}
\end{exercise}

\begin{exercise}[$\varepsilon$-$\delta$ Definition]
Using the formal definition, prove that $\lim_{x\to 1}(4x-3)=1$.
\end{exercise}

\begin{exercise}[Continuity]
Determine the values of $k$ for which
\[
  f(x)=\begin{cases} kx^2+2, & x\le 1 \\ 3x+k, & x>1 \end{cases}
\]
is continuous at $x=1$.
\end{exercise}

\begin{answerkey}
  \textbf{1a)} $11$\par
  \textbf{1b)} $1$\par
  \textbf{1c)} $2$\par
  \textbf{1d)} $\tfrac{1}{4}$\par
  \textbf{2)} $\delta=\varepsilon/4$\par
  \textbf{3)} $k=\tfrac{1}{2}$
\end{answerkey}

%%=============================================================
%%  CHAPTER 2 — DERIVATIVES
%%=============================================================
\setChapColor{Cap2Orange}{LightOrange}
\chapter{Derivatives}
\markboth{Ch.~2\ --\ Derivatives}{}

\begin{multicols}{2}
The \textbf{derivative} measures the instantaneous rate of change of a function.
Geometrically, it is the slope of the tangent line to the graph.

\begin{itemize}
  \item \topic{1}{Average vs. instantaneous rate of change}
  \item \topic{2}{Definition via limits}
  \item \topic{3}{Differentiation rules}
  \item \topic{4}{Chain rule}
  \item \topic{5}{Higher-order derivatives}
\end{itemize}

\columnbreak
\begin{didyouknow}
The derivative appears in Physics as \textbf{velocity} ($v=ds/dt$),
in Economics as \textbf{marginal cost}, and in Biology as
a population growth rate.
\end{didyouknow}
\end{multicols}

\secbar

\begin{lesson}{From Average Rate to Derivative}

\step{1}{The \textbf{average rate of change} of $f$ on $[a,a+h]$ is:}
\formula{%
  \frac{\Delta f}{\Delta x} = \frac{f(a+h)-f(a)}{h}}

\step{2}{Letting $h\to 0$, the average rate becomes the \textbf{derivative} at $a$:}
\formula[Definition of the Derivative]{%
  f'(a) = \lim_{h\to 0}\frac{f(a+h)-f(a)}{h}}

\step{3}{The result is the \textbf{slope} of the tangent line at $(a,f(a))$.}

\end{lesson}

\section{Differentiation Rules}

\begin{theorem}{Elementary Rules}
\begin{multicols}{2}
\begin{enumerate}
  \item $\dfrac{d}{dx}[c]=0$
  \item $\dfrac{d}{dx}[x^n]=nx^{n-1}$
  \item $\dfrac{d}{dx}[cf]=c f'$
  \item $\dfrac{d}{dx}[f+g]=f'+g'$
  \item $\dfrac{d}{dx}[fg]=f'g+fg'$
  \item $\dfrac{d}{dx}\!\left[\dfrac{f}{g}\right]=\dfrac{f'g-fg'}{g^2}$
\end{enumerate}
\end{multicols}
\end{theorem}

\begin{theorem}{Chain Rule}
If $y=f(u)$ and $u=g(x)$, then
\[
  \frac{dy}{dx} = \frac{dy}{du}\cdot\frac{du}{dx}
  \quad\Longleftrightarrow\quad
  [f(g(x))]' = f'(g(x))\cdot g'(x).
\]
\end{theorem}

\begin{multicols}{2}
\begin{example}{2.1}
Compute $\dfrac{d}{dx}\!\left[\sin(x^2)\right]$.

\medskip
$u=x^2$, $\frac{du}{dx}=2x$, $\frac{d}{du}[\sin u]=\cos u$

\result{$2x\cos(x^2)$}
\end{example}

\columnbreak

\begin{center}
\begin{tikzpicture}
\begin{axis}[
  width=0.95\linewidth, height=5.5cm,
  axis lines=center,
  xmin=-1.8, xmax=2.4,
  ymin=-1.2, ymax=6,
  samples=100,
  grid=both,
  grid style={RuleGray,line width=0.3pt},
  tick label style={font=\tiny},
  xlabel={$x$}, ylabel={$y$},
  xlabel style={font=\small},
  ylabel style={font=\small},
  legend style={font=\tiny,fill=white,draw=RuleGray,
                at={(0.98,0.98)},anchor=north east}]
  \addplot[Cap2Orange,line width=1.8pt,domain=-1.5:2.2] {x^2};
  \addplot[Cap1Blue,dashed,line width=1.2pt,domain=0.5:2.0] {2*(x-1)+1};
  \addplot[Cap2Orange,only marks,mark=*,mark size=2.5pt]
    coordinates {(1,1)};
  \legend{$f(x)=x^2$, tangent at $(1{,}1)$}
\end{axis}
\end{tikzpicture}
\end{center}
\end{multicols}

\begin{center}
\renewcommand{\arraystretch}{1.5}
\begin{tabular}{|>{\bfseries\color{Cap2Orange}}c|c|c|}
\hline
\rowcolor{Cap2Orange}
\color{white}\textbf{Function} & \color{white}\textbf{Derivative} &
\color{white}\textbf{Note}\\
\hline
$\sin x$    & $\cos x$               & trigonometric functions\\
$\cos x$    & $-\sin x$              & \\
$e^x$       & $e^x$                  & unique function $f=f'$\\
$\ln x$     & $1/x$                  & $x>0$\\
$x^n$       & $nx^{n-1}$             & power rule\\
$\arctan x$ & $\dfrac{1}{1+x^2}$    & inverse trig.\\
\hline
\end{tabular}
\end{center}

\section*{Chapter 2 Exercises}
\addcontentsline{toc}{section}{Exercises}

\begin{exercise}[Basic rules]
Differentiate each function:
\begin{tasks}(2)
  \task $f(x)=3x^4-5x^2+7$
  \task $g(x)=\sqrt{x}\cdot e^x$
  \task $h(x)=\dfrac{\ln x}{x^2}$
  \task $p(x)=\sin^3(2x)$
\end{tasks}
\end{exercise}

\begin{exercise}[Tangent line]
Find the equation of the tangent line to $f(x)=x^3-2x+1$ at $x=2$.
\end{exercise}

\begin{answerkey}
  \textbf{1a)} $12x^3-10x$\par
  \textbf{1b)} $\tfrac{e^x}{2\sqrt{x}}+\sqrt{x}e^x$\par
  \textbf{1c)} $\tfrac{1-2\ln x}{x^3}$\par
  \textbf{1d)} $6\sin^2(2x)\cos(2x)$\par
  \textbf{2)} $y=10x-11$
\end{answerkey}

%%=============================================================
%%  CHAPTER 3 — APPLICATIONS OF DERIVATIVES
%%=============================================================
\setChapColor{Cap3Green}{LightGreen}
\chapter{Applications of Derivatives}
\markboth{Ch.~3\ --\ Appl.\ of Derivatives}{}

\begin{multicols}{2}
In this chapter we use derivatives to \textbf{analyze} the behavior
of functions: finding maxima, minima, inflection points and sketching
graphs with precision.

\begin{itemize}
  \item \topic{1}{Increasing and decreasing intervals}
  \item \topic{2}{Local extrema — Derivative tests}
  \item \topic{3}{Concavity and inflection}
  \item \topic{4}{Optimization}
  \item \topic{5}{L'H\^opital's Rule}
\end{itemize}

\columnbreak
\begin{caution}
The maximum and minimum tests require care!
A critical point $f'(c)=0$ can be a
\textbf{maximum}, \textbf{minimum}, or
\textbf{saddle point}. Use the second derivative to confirm.
\end{caution}
\end{multicols}

\secbar

\section{Local Extrema}

\begin{theorem}{First Derivative Test}
Let $c$ be a critical point of $f$ ($f'(c)=0$):
\begin{itemize}
  \item If $f'$ changes from $+$ to $-$ at $c$ $\Rightarrow$ \textbf{local maximum}.
  \item If $f'$ changes from $-$ to $+$ at $c$ $\Rightarrow$ \textbf{local minimum}.
  \item If $f'$ does not change sign $\Rightarrow$ \textbf{saddle point}.
\end{itemize}
\end{theorem}

\begin{theorem}{Second Derivative Test}
\[
  f''(c)>0 \implies \text{local minimum};\qquad
  f''(c)<0 \implies \text{local maximum}.
\]
\end{theorem}

\begin{example}{3.1 — Optimization}
An open-top box with a square base of side $x$ and height $h$
must have volume $V=500\,\text{cm}^3$. Minimize the surface area.

\smallskip
\textbf{Solution:}\ \ $A = x^2 + 4xh$,\ $h=500/x^2$
\[
  A(x) = x^2 + \frac{2000}{x},\qquad
  A'(x) = 2x - \frac{2000}{x^2} = 0 \implies x = \sqrt[3]{1000}=10
\]

\result{$x=10\,\text{cm}$, $h=5\,\text{cm}$, $A_{\min}=300\,\text{cm}^2$}
\end{example}

\begin{center}
\begin{tikzpicture}
\begin{axis}[
  width=0.65\linewidth, height=5cm,
  axis lines=center,
  xmin=0, xmax=18, ymin=0, ymax=700,
  samples=100,
  xlabel={$x$ (cm)}, ylabel={$A(x)$},
  xlabel style={font=\small},
  ylabel style={font=\small},
  tick label style={font=\small},
  grid=both,
  grid style={RuleGray,line width=0.3pt},
  legend style={font=\small,fill=white,draw=RuleGray}]
  \addplot[Cap3Green,line width=2pt,domain=1:17] {x^2+2000/x};
  \addplot[AccentRed,dashed,line width=1.2pt,domain=0:18] {300};
  \addplot[Cap3Green,only marks,mark=*,mark size=3pt]
    coordinates {(10,300)};
  \legend{$A(x)=x^2+\frac{2000}{x}$, $A_{\min}=300$}
\end{axis}
\end{tikzpicture}
\end{center}

\section{L'H\^opital's Rule}

\begin{theorem}{L'H\^opital's Rule}
If $\lim_{x\to a}f(x)=\lim_{x\to a}g(x)=0$ (or $\pm\infty$) and
$g'(x)\neq 0$ on a neighborhood of $a$, then
\[
  \lim_{x\to a}\frac{f(x)}{g(x)} = \lim_{x\to a}\frac{f'(x)}{g'(x)}.
\]
\end{theorem}

\begin{example}{3.2}
$\displaystyle\lim_{x\to 0}\dfrac{e^x-1-x}{x^2}$
is of the form $0/0$. Applying L'H\^opital twice:
\[
  \lim_{x\to 0}\frac{e^x-1}{2x} = \lim_{x\to 0}\frac{e^x}{2} = \frac{1}{2}.
\]
\result{$1/2$}
\end{example}

\section*{Chapter 3 Exercises}
\addcontentsline{toc}{section}{Exercises}

\begin{exercise}[Extrema]
Find all local and global extrema of
$f(x)=2x^3-9x^2+12x-4$ on $[0,3]$.
\end{exercise}

\begin{exercise}[Optimization]
A farmer wants to enclose a rectangular area of $3600\,\text{m}^2$
using the minimum length of fencing.
What should the dimensions be?
\end{exercise}

\begin{exercise}[L'H\^opital]
Evaluate: (a) $\lim_{x\to 0}\dfrac{\sin 3x}{x}$;\quad
           (b) $\lim_{x\to\infty} x e^{-x}$.
\end{exercise}

\begin{answerkey}
  \textbf{1)} local min.\ $x=2$; local max.\ $x=1$;\par
  \quad global min.\ $f(0)=-4$; global max.\ $f(1)=1$\par
  \textbf{2)} $60\,\text{m}\times 60\,\text{m}$ (square)\par
  \textbf{3a)} $3$\quad\textbf{3b)} $0$
\end{answerkey}

%%=============================================================
%%  CHAPTER 4 — INDEFINITE INTEGRALS
%%=============================================================
\setChapColor{Cap4Purple}{LightPurple}
\chapter{Indefinite Integrals}
\markboth{Ch.~4\ --\ Indefinite Integrals}{}

\begin{multicols}{2}
The \textbf{indefinite integral} (or antiderivative) answers the question:
``which function, when differentiated, gives $f(x)$?''

\begin{itemize}
  \item \topic{1}{Antiderivatives and $\int$ notation}
  \item \topic{2}{Basic integrals}
  \item \topic{3}{Substitution ($u$-substitution)}
  \item \topic{4}{Integration by parts}
  \item \topic{5}{Partial fractions}
\end{itemize}

\columnbreak
\begin{didyouknow}
The \textbf{Fundamental Theorem of Calculus} links the integral to the derivative:
\[
  \frac{d}{dx}\!\int_a^x f(t)\,dt = f(x).
\]
This result is considered one of the deepest in Mathematics!
\end{didyouknow}
\end{multicols}

\secbar

\begin{lesson}{Basic Integrals}
\begin{center}
\renewcommand{\arraystretch}{1.8}
\begin{tabular}{|>{\bfseries\color{Cap4Purple}}c|c|c|}
\hline
\rowcolor{Cap4Purple}
\color{white}\textbf{Function $f(x)$} &
\color{white}\textbf{$\int f(x)\,dx$} &
\color{white}\textbf{Condition}\\
\hline
$x^n$                    & $\dfrac{x^{n+1}}{n+1}+C$ & $n\neq -1$\\
$\frac{1}{x}$            & $\ln|x|+C$               & $x\neq 0$\\
$e^x$                    & $e^x+C$                  & \\
$\sin x$                 & $-\cos x+C$              & \\
$\cos x$                 & $\sin x+C$               & \\
$\sec^2 x$               & $\tan x+C$               & \\
$\frac{1}{\sqrt{1-x^2}}$ & $\arcsin x+C$            & $|x|<1$\\
$\frac{1}{1+x^2}$        & $\arctan x+C$            & \\
\hline
\end{tabular}
\end{center}
\end{lesson}

\section{Substitution ($u$-substitution)}

\begin{theorem}{Substitution Technique}
If $u=g(x)$, then $du=g'(x)\,dx$ and
\[
  \int f(g(x))\,g'(x)\,dx = \int f(u)\,du.
\]
\end{theorem}

\begin{multicols}{2}
\begin{example}{4.1}
Evaluate $\displaystyle\int 2x\cos(x^2)\,dx$.

\textbf{Solution:}\ $u=x^2$, $du=2x\,dx$
\[
  \int\cos u\,du = \sin u + C = \sin(x^2)+C.
\]
\end{example}

\columnbreak
\begin{example}{4.2}
Evaluate $\displaystyle\int\frac{x}{x^2+1}\,dx$.

\textbf{Solution:}\ $u=x^2+1$, $du=2x\,dx$
\[
  \frac{1}{2}\int\frac{du}{u} = \frac{1}{2}\ln|u|+C
    = \frac{\ln(x^2+1)}{2}+C.
\]
\end{example}
\end{multicols}

\section{Integration by Parts}

\begin{theorem}{Integration by Parts}
\formula[Formula (product rule for integration)]{%
  \int u\,dv = uv - \int v\,du}
\textbf{LIATE Rule} for choosing $u$:
\textbf{L}ogarithmic · \textbf{I}nverse trig. · \textbf{A}lgebraic ·
\textbf{T}rigonometric · \textbf{E}xponential.
\end{theorem}

\begin{example}{4.3}
Evaluate $\displaystyle\int x e^x\,dx$.

$u=x$, $dv=e^x dx$ $\Rightarrow$ $du=dx$, $v=e^x$
\[
  \int x e^x\,dx = xe^x - \int e^x\,dx = xe^x - e^x + C = e^x(x-1)+C.
\]
\result{$e^x(x-1)+C$}
\end{example}

\section*{Chapter 4 Exercises}
\addcontentsline{toc}{section}{Exercises}

\begin{exercise}[Basic integrals]
Evaluate:
\begin{tasks}(2)
  \task $\displaystyle\int(3x^2-2x+5)\,dx$
  \task $\displaystyle\int\frac{dx}{\sqrt{4-x^2}}$
  \task $\displaystyle\int e^{3x}\,dx$
  \task $\displaystyle\int\sec^2(5x)\,dx$
\end{tasks}
\end{exercise}

\begin{exercise}[Substitution]
Evaluate $\displaystyle\int\frac{\ln x}{x}\,dx$.
\end{exercise}

\begin{exercise}[By parts]
Evaluate $\displaystyle\int x^2\ln x\,dx$.
\end{exercise}

\begin{answerkey}
  \textbf{1a)} $x^3-x^2+5x+C$\par
  \textbf{1b)} $\arcsin(x/2)+C$\par
  \textbf{1c)} $\tfrac{1}{3}e^{3x}+C$\par
  \textbf{1d)} $\tfrac{1}{5}\tan(5x)+C$\par
  \textbf{2)} $\tfrac{(\ln x)^2}{2}+C$\par
  \textbf{3)} $\tfrac{x^3}{3}\ln x - \tfrac{x^3}{9}+C$
\end{answerkey}

%%=============================================================
%%  BACK MATTER / REFERENCES
%%=============================================================
\backmatter
\setChapColor{Cap1Blue}{LightBlue}

\chapter*{References}
\addcontentsline{toc}{chapter}{References}

\begin{tcolorbox}[enhanced,colback=BoxGray,colframe=RuleGray,arc=6pt]
\begin{enumerate}[leftmargin=1.5em]
  \item STEWART, J. \textit{Calculus}. 8th ed. Cengage, 2016.
  \item ANTON, H.; BIVENS, I.; DAVIS, S. \textit{Calculus}. 10th ed.
        Wiley, 2012.
  \item APOSTOL, T. \textit{Calculus}. 2nd ed. Wiley, 1967. 2 vols.
  \item SPIVAK, M. \textit{Calculus}. 4th ed. Publish or Perish, 2008.
  \item THOMAS, G. \textit{Thomas' Calculus}. 14th ed. Pearson, 2018.
\end{enumerate}
\end{tcolorbox}

\vspace{1cm}
\begin{center}
\begin{tcolorbox}[enhanced,width=0.6\linewidth,
  colback=Cap1Blue,colframe=Cap1Blue,
  arc=12pt,halign=center,
  left=12pt,right=12pt,top=10pt,bottom=10pt]
  \color{white}\large\bfseries
  ``Calculus is the language in which\\
  God wrote the Universe.''\\[6pt]
  \small\normalfont\color{white!80!Cap1Blue}--- Galileo Galilei (attributed)
\end{tcolorbox}
\end{center}

\vspace{1cm}
\begin{center}
\qrbox{https://www.wolframalpha.com}{Wolfram Alpha ---\\
  explore functions online}
\hspace{1cm}
\qrbox{https://www.geogebra.org}{GeoGebra ---\\
  interactive visualizations}
\end{center}

\end{document}}
}